\section{Measurement Guidance for Workshop Use}

\paragraph{What this paper contributes to AIMS-style evaluation.}
The practical contribution is not only a Lean-specific score.
It is a reusable measurement workflow for evaluator governance.
Given a new judge candidate, the evaluator is not deployed directly from headline correlation.
Instead, the evaluator is screened by property tests, assigned a trust score, and compared against alternative judge configurations before being used for large-scale ranking.
This separates \emph{model quality} from \emph{measurement quality}, which is critical in settings where evaluation itself can be optimized against.

\paragraph{Validity decomposition.}
We recommend treating judge reliability as a vector rather than a scalar.
Let
\[
\mathbf{v}_\theta = (P_1, P_2, P_3, P_4, H_\theta),
\]
where \(H_\theta\) is human-alignment performance on the development panel.
Trust-index aggregation is useful for selection, but individual coordinates identify failure modes.
For example, two judges may have similar aggregate trust values but different weaknesses:
one may fail stability while another fails monotonicity under bug ladders.
These cases imply different mitigation strategies and different deployment risks.

\paragraph{A deployment protocol for strategic settings.}
In strategically optimized environments, agents can learn evaluator quirks.
We therefore recommend a staged deployment policy:
\begin{enumerate}
\item \textbf{Pre-screening.} Eliminate judge configurations that fail hard minimum thresholds on identity and monotonicity.
\item \textbf{Human-agreement ranking.} Rank remaining judges by trust index and human alignment.
\item \textbf{Shadow evaluation.} Run at least one alternate judge in parallel for disagreement monitoring.
\item \textbf{Drift alarms.} Trigger reassessment when disagreement rates or property scores shift materially.
\end{enumerate}
This policy is intentionally lightweight and can be integrated into existing benchmark pipelines with minimal annotation overhead.

\paragraph{Interactive measurement implications.}
Formal-verification pipelines are interactive because an agent's next action depends on prior tool feedback, compiler errors, and judge outputs.
A judge used inside this loop can alter the trajectory of generated artifacts.
Property-based certification is therefore not merely a post-hoc reporting tool.
It is a control mechanism for the feedback signal that drives agent behavior.
In this view, P1--P4 act as guardrails for the evaluator channel in an adaptive system.

\paragraph{Non-stationarity and reassessment cadence.}
Evaluator quality drifts for at least three reasons:
model updates, prompt updates, and benchmark mixture updates.
A simple periodic reassessment schedule can reduce silent failures.
Concretely, we recommend:
\begin{itemize}
\item fixed-interval reassessment (for example monthly or per major model release);
\item event-triggered reassessment after significant prompt or rubric edits;
\item domain-triggered reassessment when new task families are added.
\end{itemize}
The development panel can remain small if property suites are informative, but it should include anchor items that persist over time to enable longitudinal comparability.

\paragraph{Threat model and failure taxonomy.}
Certification should report which failure modes were actually tested.
For this study, the protocol directly probes:
\begin{itemize}
\item \textbf{self-consistency failure} (identity violations);
\item \textbf{directional failure} (non-monotone scoring under controlled degradation);
\item \textbf{variance failure} (instability across repeated runs).
\end{itemize}
It does not yet comprehensively test adversarial prompt attacks, contamination from known benchmark items, or distribution shifts to unseen theorem styles.
These omissions should be treated as explicit residual risk, not implicit coverage.

\paragraph{Why this is measurement science rather than only benchmarking.}
Benchmarking asks, ``which method scored highest under a chosen evaluator?''
Measurement science asks, ``why should this evaluator be trusted, under what assumptions, and how will that trust be monitored over time?''
Our position is that modern AI evaluation needs both.
Without evaluator certification, method rankings can look precise while remaining epistemically fragile.
With certification, ranking statements become conditional but auditable.

\paragraph{Practical checklist for workshop submissions.}
To make measurement claims reproducible, we suggest that judge-evaluation papers report:
\begin{enumerate}
\item the exact judge configuration tuple \((\text{model}, \text{prompt}, \text{repeats}, \text{aggregation})\);
\item property definitions and perturbation generation rules;
\item per-property scores and uncertainty summaries;
\item development-panel composition and annotation protocol;
\item drift monitoring and re-certification criteria.
\end{enumerate}
These fields are lightweight to report and substantially improve comparability across studies.

\paragraph{Scope of claims in this paper.}
Our empirical claims are intentionally bounded.
We show that a property-certified judge can track human alignment on a fixed development suite and can produce informative method rankings on \textsc{VeriBench}.
We do not claim universal validity across all formal domains or all future judge families.
The broader claim is methodological: property-based certification can serve as a robust measurement layer for AI evaluation pipelines where interaction, optimization pressure, and drift are first-order concerns.
